\documentclass[runningheads]{llncs}

% ── Packages ──────────────────────────────────────────────────────────────────
\usepackage{graphicx}
\usepackage{booktabs}
\usepackage[pagebackref,breaklinks,colorlinks,citecolor=blue]{hyperref}
\usepackage{amsmath,amssymb}
\usepackage{xcolor}
\usepackage[table]{xcolor}
\usepackage{multirow}
\usepackage{array}
\usepackage{adjustbox}
\usepackage{subcaption}
\usepackage{enumitem}

\newcommand{\tabtight}{\setlength{\tabcolsep}{3.5pt}\renewcommand{\arraystretch}{1.12}}
\newcommand{\TDS}{\textsc{TDS}}

\begin{document}

\title{Supplementary Material:\\
  Spatiotemporal Aliasing in Video Action Recognition:\\
  A Cross-Architecture Analysis at Scale}

\author{Anonymous ACCV Submission}
\institute{}
\maketitle

\setcounter{section}{0}
\renewcommand\thesection{S\arabic{section}}
\renewcommand\thefigure{S\arabic{figure}}
\renewcommand\thetable{S\arabic{table}}
\setcounter{figure}{0}
\setcounter{table}{0}

This document provides supplementary material for the main paper.
We include:
(S1) full coverage$\times$stride heatmaps for all 8 models and 7 datasets;
(S2) Levene variance inflation analysis;
(S3) full ANOVA effect size tables;
(S4) action sensitivity taxonomy detail;
(S5) entropy routing curves for all models;
(S6) clip duration analysis;
(S7) full P3 retraining results;
(S8) full E7 routing results per model/dataset;
and (S9) implementation details.

% ── S1: Full E1 heatmaps ──────────────────────────────────────────────────────
\section{Full Coverage$\times$Stride Heatmaps (\texorpdfstring{$\mathcal{E}_1$}{E1})}
\label{sec:sup_heatmaps}

Figures~\ref{fig:heatmap_r3d}--\ref{fig:heatmap_videomamba} show the complete
coverage$\times$stride accuracy heatmaps for all 8 architectures across all 7 datasets.
Each cell represents Top-1 accuracy (\%) at a specific (coverage, stride) configuration.
Color scale: green=high accuracy, red=low accuracy.
† denotes VideoMamba/AUTSL feature collapse (K400$\to$sign language domain gap).

Row order (coverage): 10\%, 25\%, 50\%, 75\%, 100\%.
Column order (stride): 1, 2, 4, 8, 16.

\begin{figure}[p]
  \centering
  \includegraphics[width=\textwidth]{images/sup1_heatmap_r3d_18.pdf}
  \caption{R3D-18 (CNN-3D) coverage$\times$stride accuracy (\%) across 7 datasets.
    Sharp cliff at stride 4$\to$8 on AUTSL (75\%$\to$27\%) confirms
    the high temporal demand of sign language recognition.
    EPIC-Kitchens and DriveAct retain high accuracy even at sparse sampling.}
  \label{fig:heatmap_r3d}
\end{figure}

\begin{figure}[p]
  \centering
  \includegraphics[width=\textwidth]{images/sup1_heatmap_mc3_18.pdf}
  \caption{MC3-18 (CNN-mix, channel-separated 3D) heatmaps.
    Similar pattern to R3D-18 but with reduced aliasing on low-TDS datasets
    (UCF-101, DriveAct), consistent with the regularizing effect of 2D-channel separation.}
  \label{fig:heatmap_mc3}
\end{figure}

\begin{figure}[p]
  \centering
  \includegraphics[width=\textwidth]{images/sup1_heatmap_r2plus1d_18.pdf}
  \caption{R2+1D-18 (separable spatiotemporal convolution) heatmaps.
    Higher base accuracy than R3D-18 on several datasets (SSv2: 42.6\% vs 37.1\%),
    but similar aliasing sensitivity pattern—cliff at stride 4$\to$8 on AUTSL.}
  \label{fig:heatmap_r2p1d}
\end{figure}

\begin{figure}[p]
  \centering
  \includegraphics[width=\textwidth]{images/sup1_heatmap_slowfast_r50.pdf}
  \caption{SlowFast-R50 (dual-pathway) heatmaps.
    Cliff occurs at stride 2$\to$4 on AUTSL (77\%$\to$41\%)—one step earlier than
    other CNNs—because the fast pathway requires 32 frames and is simultaneously
    starved along with the slow 8-frame pathway. Strong performance on Diving-48
    confirms the dual-pathway architecture captures fine-grained temporal structure
    when sampling is dense.}
  \label{fig:heatmap_slowfast}
\end{figure}

\begin{figure}[p]
  \centering
  \includegraphics[width=\textwidth]{images/sup1_heatmap_timesformer.pdf}
  \caption{TimeSformer (divided space--time attention) heatmaps.
    Notably flat across most stride values, with steep degradation only at the
    combined extreme of low coverage \emph{and} high stride.
    On Diving-48, TimeSformer maintains 36\% accuracy even at stride=16 (vs.\ 5\%
    for SlowFast), demonstrating the robustness conferred by global temporal attention.}
  \label{fig:heatmap_timesformer}
\end{figure}

\begin{figure}[p]
  \centering
  \includegraphics[width=\textwidth]{images/sup1_heatmap_vivit.pdf}
  \caption{ViViT (factorized attention) heatmaps.
    Despite being a Transformer, ViViT aliases nearly as severely as 3D CNNs
    (AUTSL: cliff at stride 2$\to$4, identical to SlowFast).
    Factorized temporal attention processes frames quasi-independently before global
    aggregation, leaving temporal reasoning vulnerable to frame dropout.}
  \label{fig:heatmap_vivit}
\end{figure}

\begin{figure}[p]
  \centering
  \includegraphics[width=\textwidth]{images/sup1_heatmap_videomae.pdf}
  \caption{VideoMAE (masked autoencoder) heatmaps.
    High base accuracy (SSv2: 52.4\%, UCF-101: 95.4\%) but steep temporal degradation
    (AUTSL cliff at stride 4$\to$8; DriveAct at stride 4$\to$8).
    The MAE pretraining objective is spatial—reconstructing patches—but does not
    prepare the model for temporal gaps, explaining the disconnect between spatial
    robustness and temporal fragility.}
  \label{fig:heatmap_videomae}
\end{figure}

\begin{figure}[p]
  \centering
  \includegraphics[width=\textwidth]{images/sup1_heatmap_videomamba.pdf}
  \caption{VideoMamba (SSM selective scan) heatmaps.
    AUTSL: feature collapse (K400 pretraining fails to transfer to sign language;
    accuracy remains at 0.4\% regardless of sampling configuration).
    On all other datasets, VideoMamba shows a flat heatmap with minimal degradation
    across strides—the SSM recurrent state aggregates temporal context continuously,
    providing natural robustness to missing frames.}
  \label{fig:heatmap_videomamba}
\end{figure}

% ── S2: Levene variance ───────────────────────────────────────────────────────
\section{Levene Variance Inflation Analysis (\texorpdfstring{$\mathcal{E}_2$}{E2})}
\label{sec:sup_levene}

Figure~\ref{fig:levene} plots inter-class standard deviation at stride=1 (x-axis)
versus stride=16 (y-axis) for all model-dataset pairs.
Points above the diagonal indicate that sparse sampling increases inter-class variability.
Significant pairs ($p < 0.05$, Bonferroni-corrected) are highlighted in red.

\begin{figure}[t]
  \centering
  \includegraphics[width=0.8\columnwidth]{images/sup2_levene_variance.pdf}
  \caption{Levene test: inter-class accuracy std at stride=1 vs.\ stride=16.
    Points above the dashed diagonal have higher variance under sparse sampling.
    Color encodes the variance ratio ($\sigma_{s16}/\sigma_{s1}$).
    \textbf{Key finding}: 67\% of model-dataset pairs show significant variance inflation
    ($p<0.05$), with the largest increases in VideoMAE/HMDB-51 (2.0$\times$) and
    R3D-18/HMDB-51 (1.85$\times$). TimeSformer and VideoMamba show stable variance
    (ratio $\approx$1.1$\times$), confirming their practical robustness.}
  \label{fig:levene}
\end{figure}

Table~\ref{tab:levene_top} lists the 10 model-dataset pairs with highest variance
inflation, quantifying the instability masked by mean accuracy.

\begin{table}[t]
\centering
\caption{Top-10 variance inflation cases (stride=1$\to$16, coverage=100\%).
  All rows significant at $p{<}0.05$.}
\label{tab:levene_top}
\scriptsize\tabtight
\begin{tabular}{l l r r r r}
\toprule
Model & Dataset & $\sigma_{s=1}$ & $\sigma_{s=16}$ & Ratio & Levene $p$ \\
\midrule
VideoMAE    & HMDB-51      & 0.125 & 0.253 & 2.02$\times$ & $<$0.0001 \\
R3D-18      & HMDB-51      & 0.164 & 0.303 & 1.85$\times$ & 0.0001 \\
ViViT       & HMDB-51      & 0.158 & 0.277 & 1.75$\times$ & 0.0004 \\
SlowFast    & HMDB-51      & 0.187 & 0.326 & 1.74$\times$ & $<$0.0001 \\
SlowFast    & UCF-101      & 0.146 & 0.240 & 1.65$\times$ & $<$0.0001 \\
R2+1D       & HMDB-51      & 0.205 & 0.310 & 1.51$\times$ & 0.0021 \\
MC3-18      & HMDB-51      & 0.177 & 0.262 & 1.48$\times$ & 0.0187 \\
ViViT       & UCF-101      & 0.110 & 0.147 & 1.34$\times$ & 0.0093 \\
VideoMAE    & AUTSL        & 0.172 & 0.229 & 1.33$\times$ & 0.0324 \\
VideoMAE    & UCF-101      & 0.094 & 0.121 & 1.29$\times$ & 0.0239 \\
\bottomrule
\end{tabular}
\end{table}

% ── S3: Full ANOVA tables ─────────────────────────────────────────────────────
\section{Full ANOVA Effect Size Tables (\texorpdfstring{$\mathcal{E}_4$}{E4})}
\label{sec:sup_anova}

Figure~\ref{fig:anova_bar} shows the mean $\eta^2$ per model (averaged over all valid
datasets), with error bars indicating standard deviation.

\begin{figure}[t]
  \centering
  \includegraphics[width=0.9\columnwidth]{images/sup3_anova_eta2.pdf}
  \caption{ANOVA effect sizes: $\eta^2(\text{coverage})$ + $\eta^2(\text{stride})$
    per model. Coverage consistently dominates (blue bars, 53--90\% of variance).
    The stride effect is large for CNNs and ViViT (red bars, 18--35\%) but small
    for VideoMamba and TimeSformer ($<$10\%), confirming architecture-dependent aliasing.
    Error bars: std across 7 datasets.}
  \label{fig:anova_bar}
\end{figure}

Table~\ref{tab:anova_full} provides the complete per-dataset $\eta^2$ values.
All stride effects are significant ($p < 0.05$); coverage effects are significant
($p < 0.001$) in all cases.

\begin{table}[t]
\centering
\caption{Full ANOVA $\eta^2$ per model and dataset. Cov = $\eta^2(\text{coverage})$,
  Str = $\eta^2(\text{stride})$. Rows shaded darker indicate larger stride effect.}
\label{tab:anova_full}
\scriptsize\tabtight
\begin{adjustbox}{max width=\textwidth}
\begin{tabular}{l *{7}{cc}}
\toprule
\multirow{2}{*}{Model}
 & \multicolumn{2}{c}{AUTSL}
 & \multicolumn{2}{c}{Diving-48}
 & \multicolumn{2}{c}{SSv2}
 & \multicolumn{2}{c}{HMDB-51}
 & \multicolumn{2}{c}{DriveAct}
 & \multicolumn{2}{c}{EPIC}
 & \multicolumn{2}{c}{UCF-101} \\
\cmidrule(lr){2-3}\cmidrule(lr){4-5}\cmidrule(lr){6-7}\cmidrule(lr){8-9}
\cmidrule(lr){10-11}\cmidrule(lr){12-13}\cmidrule(lr){14-15}
 & Cov & Str & Cov & Str & Cov & Str & Cov & Str & Cov & Str & Cov & Str & Cov & Str \\
\midrule
VMamba & .88 & .09 & .94 & .04 & .88 & .09 & .90 & .09 & .83 & .15 & .95 & .05 & .88 & .07 \\
TSF    & .92 & .06 & .97 & .02 & .87 & .10 & .92 & .07 & .84 & .14 & .91 & .08 & .85 & .07 \\
MC3    & .66 & .18 & .81 & .13 & .54 & .28 & .81 & .18 & .75 & .22 & .89 & .11 & .79 & .18 \\
R3D    & .60 & .21 & .77 & .15 & .60 & .26 & .68 & .31 & .58 & .33 & .75 & .24 & .70 & .25 \\
R2+1D  & .63 & .20 & .76 & .17 & .62 & .26 & .72 & .27 & .49 & .38 & .80 & .18 & .76 & .20 \\
VMAE   & .69 & .18 & .82 & .12 & .69 & .22 & .84 & .16 & .43 & .43 & .87 & .11 & .81 & .14 \\
ViViT  & .61 & .24 & .69 & .20 & .51 & .33 & .71 & .27 & .60 & .31 & .80 & .19 & .68 & .26 \\
SF     & .51 & .29 & .51 & .33 & .44 & .36 & .55 & .39 & .39 & .46 & .67 & .30 & .61 & .36 \\
\bottomrule
\end{tabular}
\end{adjustbox}
\end{table}

% ── S4: Taxonomy ──────────────────────────────────────────────────────────────
\section{Action Sensitivity Taxonomy (\texorpdfstring{$\mathcal{E}_5$}{E5})}
\label{sec:sup_taxonomy}

\begin{figure}[t]
  \centering
  \includegraphics[width=\textwidth]{images/sup4_taxonomy.pdf}
  \caption{Aliasing loss (pp) per sensitivity tier and dataset.
    Each bar shows mean accuracy drop (stride=1$\to$16) within the tier,
    annotated with the number of action classes.
    UCF-101 Low tier: $-0.3$pp (static, appearance-dominated actions lose
    essentially no accuracy under sparse sampling).
    AUTSL: all three tiers remain high ($>38$pp), confirming that sign language
    is fundamentally high-frequency regardless of action subclass.}
  \label{fig:taxonomy_full}
\end{figure}

Table~\ref{tab:taxonomy_detail} lists the threshold values that define tier boundaries
per dataset, derived from the 33rd and 67th percentiles of the per-class accuracy drop.

\begin{table}[t]
\centering
\caption{Taxonomy tier thresholds (accuracy drop, pp): classes below the Low threshold
  alias negligibly; classes above the High threshold are severely affected.
  Thresholds are dataset-specific and computed across all 8 valid architectures.}
\label{tab:taxonomy_detail}
\scriptsize\tabtight
\begin{tabular}{l r r r r r}
\toprule
Dataset & Classes & Low threshold & High threshold & High count & Low count \\
\midrule
AUTSL         & 226 & 66.9pp & 73.9pp & 75 & 74 \\
Diving-48     &  47 & 25.0pp & 47.7pp & 16 & 16 \\
SSv2          & 174 & 59.7pp & 76.9pp & 58 & 57 \\
HMDB-51       &  51 &  8.4pp & 28.5pp & 17 & 17 \\
EPIC-Kitchens &  89 &  0.0pp & 19.4pp & 30 & 11 \\
DriveAct      &  33 & 14.8pp & 51.0pp & 11 & 11 \\
UCF-101       & 101 &  1.4pp &  7.5pp & 35 & 31 \\
\bottomrule
\end{tabular}
\end{table}

% ── S5: E7 routing curves ─────────────────────────────────────────────────────
\section{Entropy Routing Curves — All Models (\texorpdfstring{$\mathcal{E}_7$}{E7})}
\label{sec:sup_routing}

Figure~\ref{fig:routing_all} shows accuracy vs.\ average frames for all 8 architectures
on SSv2 under the entropy routing protocol.
Each curve is the result of sweeping confidence threshold $\tau \in [0.10, 0.95]$.
Dashed horizontal lines indicate fixed-budget baselines (4f, 16f) and the oracle.

\begin{figure}[t]
  \centering
  \includegraphics[width=\textwidth]{images/sup5_routing_all_models.pdf}
  \caption{Entropy routing: accuracy vs.\ avg frames (SSv2, 4f cheap $\to$ 16f dense).
    For architectures with flat 4f$\approx$16f accuracy (TimeSformer, VideoMamba),
    routing near-trivially selects 4f for most videos and matches dense accuracy.
    For architectures with large 4f--16f gaps (SlowFast, ViViT), the routing curve
    provides meaningful operating points between the fixed-budget extremes.}
  \label{fig:routing_all}
\end{figure}

Table~\ref{tab:routing_full} provides the full E7 routing results across all models
and datasets at the optimal threshold (avg frames $\leq 8$).

\begin{table}[t]
\centering
\caption{E7 entropy routing: best accuracy at avg $\leq8$ frames per model/dataset.
  Values are Top-1 accuracy (\%). Bold = highest per column. † AUTSL excluded
  for VideoMamba (feature collapse).}
\label{tab:routing_full}
\scriptsize\tabtight
\begin{adjustbox}{max width=\textwidth}
\begin{tabular}{l ccccccc r}
\toprule
Model & AUTSL & Diving & SSv2 & HMDB & DriveAct & EPIC & UCF & Avg \\
\midrule
R3D-18       & 73.3 & 27.6 & 34.6 & 80.2 & 66.1 & 36.7 & 81.6 & 57.2 \\
MC3-18       & 63.4 & 31.8 & 32.4 & 78.5 & 69.4 & 36.5 & 85.6 & 56.8 \\
R2+1D        & 74.2 & 36.3 & 40.8 & 73.3 & 59.6 & 34.8 & 89.0 & 58.3 \\
SlowFast     & 59.5 & 46.6 & 35.0 & 76.3 & 67.4 & 37.4 & 87.9 & 58.6 \\
TimeSformer  & 67.1 & 38.9 & 42.5 & 80.6 & 68.5 & 32.4 & 91.3 & \textbf{60.2} \\
ViViT        & 67.2 & 48.9 & 32.2 & 80.1 & 64.3 & 31.3 & 94.3 & 59.8 \\
VideoMAE     & \textbf{79.5} & \textbf{52.6} & \textbf{51.9} & \textbf{84.7} & \textbf{69.0} & \textbf{37.8} & \textbf{95.9} & \textbf{67.3} \\
VideoMamba   & †    & 37.4 & 44.4 & 69.8 & 57.6 & 28.4 & 88.5 & 54.4 \\
\bottomrule
\end{tabular}
\end{adjustbox}
\end{table}

% ── S6: Clip duration ─────────────────────────────────────────────────────────
\section{Clip Duration Analysis (\texorpdfstring{$\mathcal{E}_{10}$}{E10})}
\label{sec:sup_duration}

\begin{figure}[t]
  \centering
  \includegraphics[width=\textwidth]{images/sup6_clip_duration.pdf}
  \caption{Aliasing loss (pp) vs.\ clip duration bin for all 8 architectures across
    multiple datasets. Counter-intuitive finding: \textbf{shorter clips alias more}
    (Pearson $r=-0.3$ to $-0.8$ across models). Short clips have less temporal
    redundancy; any frame drop removes a larger fraction of the available evidence.
    Long clips (${>}4$s) contain natural temporal repetition that compensates for
    sparse sampling. Exception: Diving-48 longer clips ($>$4s) alias more,
    consistent with longer dives containing more distinct phase transitions.}
  \label{fig:duration}
\end{figure}

% ── S7: P3 full results ───────────────────────────────────────────────────────
\section{Resolution Retraining Results (P3)}
\label{sec:sup_p3}

Table~\ref{tab:p3_full} provides all available P3 checkpoint results (current as of
paper submission; ongoing retraining in progress for Transformers and VideoMamba).
Each entry shows the best validation accuracy from 10 epochs of training.
``Native'' accuracy is from the E1 sweep at stride=1, coverage=100\%.

\textbf{Key patterns:}
\begin{itemize}[itemsep=2pt]
\item \textbf{96px consistently beats native for CNNs} across most datasets:
  regularization effect is strongest on background-noisy datasets (EPIC-Kitchens,
  DriveAct) and weaker on datasets requiring fine spatial detail (AUTSL, HMDB-51).
\item \textbf{160px also beats native} for R3D-18 and MC3-18, suggesting the
  regularization effect persists across a range of sub-native resolutions.
\item \textbf{AUTSL is the exception}: sign language requires fine handshape
  discrimination; very small resolutions ($<$96px) would likely degrade further.
\end{itemize}

\begin{table}[t]
\centering
\caption{P3 retraining: best Top-1 val accuracy (\%) at each resolution.
  Native = accuracy from E1 (stride=1, cov=100\%). Bold = beats native.
  Entries marked ``—'' not yet completed.}
\label{tab:p3_full}
\scriptsize\tabtight
\begin{adjustbox}{max width=\textwidth}
\begin{tabular}{ll r rrrrr}
\toprule
\multirow{2}{*}{Model} & \multirow{2}{*}{Dataset}
  & \multirow{2}{*}{Native}
  & \multicolumn{4}{c}{Retrained at resolution}\\
\cmidrule(lr){4-7}
  & & & 96px & 112px & 160px & 224px & 336px \\
\midrule
\multirow{7}{*}{R3D-18}
  & UCF-101      & 81.2\% & \textbf{86.5\%} & — & \textbf{87.4\%} & — & — \\
  & SSv2         & 37.1\% & \textbf{45.2\%} & — & \textbf{45.8\%} & — & — \\
  & HMDB-51      & 80.3\% & 78.3\%           & — & 78.7\%           & — & — \\
  & Diving-48    & 28.8\% & \textbf{31.0\%} & — & \textbf{37.2\%} & — & — \\
  & AUTSL        & 75.0\% & 69.5\%           & — & \textbf{82.7\%} & — & — \\
  & DriveAct     & 68.3\% & \textbf{80.5\%} & — & \textbf{83.2\%} & — & — \\
  & EPIC-Kitchen & 37.2\% & \textbf{58.0\%} & — & \textbf{59.5\%} & — & — \\
\midrule
\multirow{7}{*}{MC3-18}
  & UCF-101      & 85.4\% & 83.9\%           & — & 84.5\%           & — & — \\
  & SSv2         & 33.6\% & \textbf{42.4\%} & — & \textbf{42.8\%} & — & — \\
  & HMDB-51      & 78.6\% & 76.9\%           & — & 76.1\%           & — & — \\
  & Diving-48    & 31.6\% & \textbf{40.0\%} & — & \textbf{45.7\%} & — & — \\
  & AUTSL        & 63.7\% & 53.3\%           & — & \textbf{69.6\%} & — & — \\
  & DriveAct     & 69.0\% & \textbf{74.9\%} & — & \textbf{79.2\%} & — & — \\
  & EPIC-Kitchen & 36.2\% & \textbf{50.8\%} & — & 46.2\%           & — & — \\
\midrule
\multirow{3}{*}{R2+1D}
  & UCF-101      & 88.6\% & \textbf{89.8\%} & — & — & — & — \\
  & SSv2         & 42.6\% & \textbf{49.6\%} & — & — & — & — \\
  & HMDB-51      & 73.1\% & \textbf{84.0\%} & — & — & — & — \\
\midrule
\multirow{7}{*}{SlowFast}
  & UCF-101      & 87.6\% & \textbf{93.5\%} & \textbf{93.6\%} & \textbf{91.7\%} & — & — \\
  & SSv2         & 49.5\% & \textbf{58.4\%} & —              & — & — & — \\
  & HMDB-51      & 79.3\% & \textbf{85.8\%} & \textbf{71.0\%} & — & — & — \\
  & Diving-48    & 50.5\% & \textbf{59.0\%} & \textbf{35.5\%} & — & — & — \\
  & AUTSL        & 82.3\% & \textbf{82.8\%} & 63.1\%          & \textbf{84.0\%} & — & — \\
  & DriveAct     & 72.5\% & \textbf{82.7\%} & \textbf{76.7\%} & — & — & — \\
  & EPIC-Kitchen & 39.4\% & \textbf{62.6\%} & —              & — & — & — \\
\midrule
VideoMAE & UCF-101 & 95.4\% & — & — & — & — & 90.6\% \\
\midrule
Transformers & All & — & \multicolumn{4}{c}{\textit{In progress (training ongoing)}} & \\
VideoMamba   & All & — & \multicolumn{4}{c}{\textit{In progress (training ongoing)}} & \\
\bottomrule
\end{tabular}
\end{adjustbox}
\end{table}

% ── S8: Spectral correlation ──────────────────────────────────────────────────
\section{Spectral Correlation Detail (\texorpdfstring{$\mathcal{E}_3$}{E3})}
\label{sec:sup_spectral}

We estimated per-class temporal frequency demand using optical-flow magnitude
(Farneback algorithm~\cite{farneback2003two}) computed on 5 representative videos
per class, sampling 16 frames per video at 112$\times$112px.
The mean inter-frame flow magnitude serves as a proxy for temporal activity level.

Table~\ref{tab:spectral} reports Pearson correlation between per-class mean flow
magnitude and aliasing loss (accuracy drop stride=1$\to$16) per dataset.
Prior studies on small closed-vocabulary datasets (UCF-101, Kinetics-400) have reported
near-perfect optical-flow--aliasing correlations ($r{=}0.99$) using controlled synthetic flow protocols.
Our large-scale evaluation on 7 diverse datasets yields moderate correlations ($r=0.03$--$0.33$).
This discrepancy has three plausible explanations:
(1) complex background motion in egocentric datasets (EPIC-Kitchens) decorrelates
flow from action content;
(2) sign language (AUTSL) involves localized, subtle movements that optical flow
at 112px partially misses;
(3) our use of raw flow magnitude rather than class-conditioned frequency
analysis is a coarser proxy.

\begin{table}[t]
\centering
\caption{Spectral correlation: Pearson $r$ between optical-flow magnitude and
per-class aliasing loss. Positive $r$ supports the Nyquist hypothesis; near-zero
values suggest flow magnitude is an insufficient proxy for temporal demand on
complex datasets.}
\label{tab:spectral}
\scriptsize\tabtight
\begin{tabular}{l r r r}
\toprule
Dataset & Pearson $r$ & $p$-value & Significant ($p{<}0.05$) \\
\midrule
DriveAct      & 0.327 & 0.065 & Marginal \\
AUTSL         & 0.285 & $<$0.001 & Yes \\
SSv2          & 0.184 & $<$0.001 & Yes \\
HMDB-51       & 0.180 & 0.206 & No \\
Diving-48     & 0.140 & 0.349 & No \\
UCF-101       & 0.031 & 0.762 & No \\
EPIC-Kitchens & --0.002 & 0.985 & No \\
\bottomrule
\end{tabular}
\end{table}

% ── S9: Implementation details ────────────────────────────────────────────────
\section{Implementation Details}
\label{sec:sup_impl}

\paragraph{Hardware.}
All training experiments were conducted on an HPC cluster with NVIDIA A100 (40GB)
and H200 NVL (141GB) GPUs using CUDA 13.2.
E1 evaluation (1,400 configurations) was parallelized across 8 nodes using a
custom Slurm auto-submitter.

\paragraph{Training hyperparameters.}
All models were fine-tuned for 10 epochs.
Transformers and VideoMamba: AdamW, lr=$2\times10^{-5}$, weight decay=0.05,
batch size=32 per GPU.
CNNs: AdamW, lr=$10^{-4}$, batch size=128 per GPU.
SlowFast: AdamW, lr=$10^{-4}$, batch size=32 per GPU.
All models use cosine learning rate decay with linear warmup (1 epoch).
Spatial augmentation: random resized crop, horizontal flip, normalization
(ImageNet statistics).

\paragraph{VideoMamba positional embedding interpolation.}
VideoMamba's positional embeddings are hardcoded for 224px input (14$\times$14=196 tokens + 1 cls).
For non-native resolutions, we load the pretrained checkpoint at 224px, then bicubically
interpolate the spatial positional embeddings from 14$\times$14 to $\lfloor r/16\rfloor^2$
before copying weights to a model instantiated at the target resolution.

\paragraph{SlowFast spatial pooling fix.}
The original SlowFast-R50 head uses a fixed $7\times7$ average pooling kernel,
assuming 224px input.
We replace both head pooling layers with \texttt{AdaptiveAvgPool3d(1)},
enabling forward passes at arbitrary input resolutions.

\paragraph{Code availability.}
Code and evaluation scripts will be released at [anonymized for review].

\bibliographystyle{splncs04}
\bibliography{main}

\end{document}
